\chapter*{Kurzfassung}
\addcontentsline{toc}{chapter}{Kurzfassung}
Die vorliegende Diplomarbeit behandelt die Konzeption und Umsetzung der Webanwendung \glqq Mediahub\grqq{}, deren Ziel in der zentralen Verwaltung persönlicher Streaming-Abonnements mit Fokus auf Filme liegt. Ausgangspunkt ist die Beobachtung, dass bei parallel genutzten Streaming-Diensten die Übersicht über Kosten, Laufzeiten, Kündigungszeitpunkte und tatsächliche Verfügbarkeiten rasch verloren geht. Daraus ergibt sich ein praktisches Problem, das sowohl organisatorische als auch technische Anforderungen umfasst.

Im Rahmen des Projekts wurde eine Anwendung entwickelt, die Filminformationen, Anbieter-Verfügbarkeiten und persönliche Abo-Daten in einer gemeinsamen Oberfläche zusammenführt. Der Schwerpunkt meines Projektteils lag auf der Frontend-Entwicklung mit Angular. Umgesetzt wurden insbesondere die Abo-Verwaltung, die Filmsuche, die Film-Detailansicht, die Darstellung passender Anbieter sowie geschützte Benutzerbereiche auf Basis von Keycloak. Dadurch entstand eine funktionale Webanwendung, die einen konkreten Mehrwert für Nutzende bietet.

Die Arbeit zeigt zugleich, dass die Entwicklung einer solchen Anwendung nicht nur aus der Implementierung einzelner Funktionen besteht. Wesentliche Herausforderungen lagen in der sinnvollen Übersetzung von Backend-Daten in die Benutzeroberfläche, in der Integration von Authentifizierung sowie in der bewussten Begrenzung des Projektumfangs. Nicht alle begonnenen Funktionen wurden vollständig ausgebaut; die Verlaufsfunktion ist beispielsweise nur teilweise umgesetzt. Insgesamt stellt die Arbeit jedoch eine fachlich nachvollziehbare und technisch tragfähige Grundlage für spätere Erweiterungen dar.

